%% kapa: exam=1 round=21 year=2010 period=1 form=A subject=civil_law subject_ko=민법 seq=14 number=14
%% source: Q-Net 감정평가사 기출문제 (cst00309) artlSeq=5024181 file=extracted/1008/5024181/2010년 제21회 감정평가사 1차시험 기출문제/1교시.hwp
%% conversion: kordoc 4.24.51 -> tools/gen1.py (auto); review=proofread-sonnet; flags=none
\begin{question}{14}
민법상 주소에 관한 설명 중 옳은 것을 모두 고른 것은?
\begin{framedtext}
ㄱ. 주소는 정주(定住)의 의사를 요건으로 한다. \par
ㄴ. 주소는 부재와 실종의 표준이 된다. \par
ㄷ. 법인의 주소는 그 주된 사무소의 소재지에 있는 것으로 한다. \par
ㄹ. 거래안전을 위해 주소는 동시에 두 곳 이상 둘 수 없다. \par
\end{framedtext}
\begin{choices}
  \item ㄱ, ㄴ
  \item ㄱ, ㄷ
  \item ㄴ, ㄷ
  \item ㄴ, ㄹ
  \item ㄷ, ㄹ
\end{choices}
\end{question}
\answer{3}
